\section{Actual form and signature spaces, with the initial theta fields explicit}
\label{sec:automorphic-realization}

\status{MR1--MR11 construct the actual definite adelic realization,
including its full range, prime-local integral Hecke action and the
pointwise \texttt{MtAd} comparison on a supplied compatible geometric
state. MR12--MR22 construct the integral form and signature portion of
\texttt{SideB} over the already chosen coefficient rings. The source's
\texttt{mkF/mkEx} initializes separate theta fields as zero and still
requires a theta witness; MR22 states this boundary literally.
MR23--MR26 prove four local Hecke-algebra and relation-lattice
primitives on the original objects. The good-row spectral input,
automorphic Galois input and compatible geometric state remain
explicit. No compatible nonempty theta/local realization or global
branch conclusion follows from these initial fields.}

Write $L=D.\mathrm{what}_2.\mathrm{LL}$, $K=D.\mathrm{LL0}\subset L$,
and $\epsilon:L\hookrightarrow\mathbb C$ for the fixed coefficient
embedding. A rational form datum has a number field $K_s\subset K$,
an $L$-space $W_L$ and a $K_s$-form $W_{K_s}\subset W_L$; ``form'' here
includes the assertion $L\otimes_{K_s}W_{K_s}\simeq W_L$.
The letter $E$ denotes the host CM field and $F=E^+$.
All the old choices of host, rows, $L$, $K$, embeddings, tuple maps and
level characters remain fixed unless an explicitly specified
\texttt{setTw} operation is made.

We use the preceding finite-class and coefficient constructions
CR1--CR6, FC/TX, AR5--AR9 and AR15--AR18 where their hypotheses hold.
For coefficient representations we retain the standard characteristic
zero highest-weight theorem in its precise polynomial form: the
Schur module of a dominant weight is irreducible, its indicated
highest-weight line is one-dimensional, and its dual has the opposite
contragredient highest weight. Tensor products at the distinct,
nonconjugate embeddings of a CM type have the corresponding product
property. This is a representation-theoretic input, not an assertion
about the existence of automorphic forms or Galois representations.
The literal source interfaces are indexed in the accompanying ledger.

\paragraph{MR1: finite class values reconstruct all adelic forms.}
Let $G$ be the finite adelic group of a positive definite unitary
space, $\Gamma$ its rational group, $J$ a compact open level, and
$\chi:J\to L^\times$ the actual finite-order level character. Fix
representatives $y_c$ of the finite set $\Gamma\backslash G/J$.
Let $V_L$ be the coefficient module, with complex realization $V$ and
rational action $\rho$. For each $c$, let $V_{L,c}$ consist of the
vectors satisfying
\[
 \rho(\gamma)v=\chi(k)v
 \quad\text{whenever }\gamma y_c=y_c k,
 \qquad \mathcal C_L=\bigoplus_c V_{L,c}.             \tag{MR.1}
\]
The sign in (MR.1) follows from the convention
$F(\gamma gk)=\rho(\gamma)F(g)\chi(k)$ and the displayed
stabilizer equation.
For $v=(v_c)$ define
\[
       \operatorname{Ext}(v)(\gamma y_c k)
          =\rho(\gamma)\epsilon(\chi(k))v_c.          \tag{MR.2}
\]
If two decompositions are used, their quotient is one of the
stabilizer pairs in (MR.1), so their values agree. Conversely, every
equivariant adelic function has precisely these stabilizer relations
at $y_c$, and (MR.2) recovers it. Evaluation and extension are inverse.
The same argument over $\mathbb C$ proves surjectivity onto the full
adelic form space, not merely onto the span of a selected collection
of forms. Local constancy follows from the open kernel of $\chi$.

The rational stabilizers are finite: they preserve a lattice, and
their matrices form a discrete subset of the product of the compact
unitary groups at all real places. Their fixed-vector equations are therefore finite linear
equations over $L$. Row reduction, or SE2, shows that their kernels
commute with $L\hookrightarrow\mathbb C$. Thus a basis of
$\mathcal C_L$ extends to a complex basis of all adelic forms.
With the actual coefficient lattices $\nabla(\Lambda_c)$, integrality
is exactly $v_c\in\nabla(\Lambda_c)$ at every class. This supplies the
evaluation, extension and scalar-extension facts used from
\texttt{ClassPres}; finiteness and the specified geometric and
coefficient-lattice presentation are still required inputs.

\paragraph{MR2: the finite level projector and its rational basis.}
Suppose the actual larger level acts on $\mathcal C_L$ through a
finite group $G_0$, with representation $a$, and let $P_0\le G_0$.
For a character $\psi:P_0\to L^\times$, put
\[
 e_\psi=\frac1{|P_0|}\sum_{p\in P_0}\psi(p)^{-1}a(p).
                                                               \tag{MR.3}
\]
Changing the variable $p$ to $q^{-1}p$ gives
$a(q)e_\psi=\psi(q)e_\psi$. If $a(q)v=\psi(q)v$ for all $q$, then
$e_\psi v=v$. Hence $e_\psi^2=e_\psi$ and its image is exactly the
specified character subspace. Choose independent columns of its
matrix and assemble them into $E_\psi:L^d\to\mathcal C_L$.
Gaussian elimination gives a left inverse over $L$ and proves that
the columns are a basis of its image.

Apply \texttt{Ext} to these columns. They remain independent after
applying $\epsilon$, because the left inverse also extends to
$\mathbb C$. Every complex form satisfying the extra level character
is fixed by the corresponding averaging operator. Average its
expansion in the full basis from MR1; the result belongs to the span
of the columns of $E_\psi$. Thus these columns span the entire cut
space. If its values at the $y_c$ belong to the rational coefficient
module, MR1 first descends the form to $\mathcal C_L$, and (MR.3)
then descends it to $E_\psi(L^d)$. Finally,
\[
 x\in\mathcal N_\psi
 \quad\Longleftrightarrow\quad
       (E_\psi x)_c\in\nabla(\Lambda_c)\quad\text{for every }c
                                                               \tag{MR.4}
\]
is the actual level lattice. No division by $|P_0|$ is asserted to
preserve this lattice; the lattice is the inverse image in (MR.4).

\paragraph{MR3: the standard-frame type space and the inverse character.}
For the source's identity Gram matrix and a unit frame $P$, changing
the tensor order gives an isomorphism $t_P:G^{\rm std}(\mathbb A_f)
\simeq G$. On rational points it is the same matrix over $E$:
injectivity of $E\to E\otimes_F\mathbb A_{F,f}$ proves both
directions of the rational-point assertion. There is therefore a
compatible rational-group isomorphism, and no approximation of a
rational element is needed.

Let $Q_0$ be the larger parabolic level, let
$r:Q_0\twoheadrightarrow G_0$ be the prescribed reduction, and let
$e_G:G_0\simeq G_0(P_g,r)$ be its actual identification with the
product residue group. The standard level is precisely
\[
 J_{\rm std}=t_P^{-1}\{q\in Q_0:r(q)\in P_0\},\qquad
 \chi_{\rm std}(g)=\operatorname{ch}(t_Pg)
                       \epsilon(\psi_0(r(t_Pg))^{-1}).\tag{MR.5}
\]
Indeed, membership says that the reduced matrix preserves the
specified flags and that its upper-left blocks lie in $e_G(P_0)$.
The source uses the conjugate level character in \texttt{levFn},
which is the inverse in (MR.5). Substitution in the two equivariance
laws now proves, in both directions,
\[
 F\in\operatorname{cutSp}(\psi_0^{-1})
 \quad\Longleftrightarrow\quad
 F\circ t_P\in\operatorname{TypeSp}
                      (\mathrm{wtRep},J_{\rm std},\chi_{\rm std}).
                                                               \tag{MR.6}
\]
The smaller original level lies in $\ker r$, so its character is
also recovered by (MR.5). Thus no additional equivariance condition
is lost in the reverse implication. MR2 is applied with
$\psi=\psi_0^{-1}$, exactly as in \texttt{SideIn.fn}.

\paragraph{MR4: genuine product levels and the local descriptions.}
The standard level in MR3 is bounded by the integral maximal compact.
Take $N$ to be the product of the norms of all primes in $P_g\cup Q$.
It is positive and belongs to each of those prime ideals. A matrix
congruent to $1$ modulo $N$ reduces to $1$ at every level prime.
It therefore belongs to $J_{\rm std}$ and its character is $1$.
This proves the literal open-level condition. At a real-field prime
below none of $P_g\cup Q$, the local group is the full integral
unitary group and its level character is trivial; this is
\texttt{HypOutside}.

For completeness, the product assertion can be proved inside the
adelic tensor algebra. Let $e_v$ be $1$ at the place $v$ and $0$
elsewhere, and put
\[
             g_v=e_vg+(1-e_v)1.                      \tag{MR.7}
\]
Since $e_v^2=e_v=\bar e_v$, the unitary equation for $g$ implies
the unitary equation for $g_v$; its inverse is obtained from $g^{-1}$
by the same formula. Both integrality and each flag condition are
tested componentwise. The two complementary cuts multiply to $g$,
and their characters multiply to the character of $g$.
At a conductor prime the reduction condition is the actual
$(r,n-r)$ parabolic, with character the prescribed component on the
determinant of the complementary block. At a type prime the
upper-left block is the specified unipotent group, and the
superdiagonal sum is read by the actual nontrivial residue additive
character. Its completion character is obtained by extending the
residue character from local integers, as proved in AR6; AR7--AR8
give the exact tuple component and sign. AR5 gives the conductor
character. These identifications require the literal half-selection,
split-place and generic-character hypotheses of \texttt{LevCh.Is};
productness alone does not supply them.

\paragraph{MR5: the actual algebraic coefficient representation.}
Let $c_0:VV_L(\lambda)\to V$ be the given semilinear coefficient
map, with complex-independent basis images spanning $V$ and with
the actual rational-action identity. Define
$T:VV_{\mathbb C}(\lambda)\simeq V$ by sending the standard basis
to those images. On this basis, and hence everywhere,
\[
       T\,\rho_{VV}(g)=\rho_V(g)T.                   \tag{MR.8}
\]
The inverse-transpose matrix convention is preserved: scalar
extension commutes with matrix inverse because the two inverse
identities remain true after applying a field homomorphism.
The polynomial matrices of $VV$ therefore give polynomial matrices
of $\rho_V$ in the transported basis. Irreducibility and the dual
highest-weight statement follow from the precise highest-weight
input stated above, transported through $T$ and the nonzero dual
pairing. Distinct nonconjugate embeddings are essential here.

The passage from rational unitary points to the product of general
linear groups has an elementary density justification. On the
rational vector space of skew-Hermitian matrices the Cayley map
$X\mapsto(1-X)(1+X)^{-1}$ is inverse to the same formula on the
open subset where $1+g$ is invertible. A polynomial vanishing on
all rational points of a vector space over $\mathbb Q$ is zero,
by induction in its coordinates. Clearing the Cayley denominators
therefore proves Zariski density of the rational unitary points in
the restriction-of-scalars group. After algebraic scalar extension
this group is a product of general linear groups, one for each
conjugate embedding pair. Projection to distinct, nonconjugate
embeddings is surjective. An invariant subspace is defined by
polynomial matrix equations in a basis and a quotient map; density
makes it invariant under that entire product. The retained
polynomial highest-weight theorem consequently applies to the
actual rational-group representation as well.

For the actual row $n=R.n_V$ and
$\lambda(\tau)=\mathrm{lamLdef}(n,j,\tau)$ the dual highest
weight is exactly $\lambda(\tau)$; substituting its entries in
$\operatorname{infOf}$ gives
\[
 \operatorname{infOf}(\lambda(\tau))
       =\operatorname{paramV}(R,h_c(\tau),j-1,1/2).   \tag{MR.9}
\]
Thus the weight in \texttt{MtReal} is the coefficient representation
already used in the geometric forms. It is not a newly chosen
irreducible representation with an equal dimension.

\paragraph{MR6: rational Hecke operators with their actual normalization.}
At a split prime outside the level, the Hecke double coset of the
specified diagonal matrix is a finite union of right $J_{\rm std}$
cosets. To see finiteness, $J\cap aJa^{-1}$ is open in compact $J$,
so its cosets have a finite subcover. Choose the actual representatives
$a_t$ supported at that prime, where the level character is trivial.
They give the source operator
\[
       T_{s,k}F(g)=q_s^{-k(n-1)}\sum_t F(ga_t).        \tag{MR.10}
\]
Changing a representative by a local level element has no effect,
since its character is trivial. Left rational equivariance commutes
with the sum. For right equivariance, the local component permutes
these cosets; the complementary level component commutes with
their supported representatives and contributes exactly its original
character. MR4 supplies this product decomposition. Thus
$T_{s,k}$ preserves the full type space.
Its normalization in (MR.10) is the literal \texttt{NG.tN}, rather
than an unnormalized double-coset sum.

Every value in the sum is in the image of the rational coefficient
module: MR1 and the actual coefficient map give this at an arbitrary
adelic argument, not just at the selected representatives.
The normalization is rational. By MR2 and type-space preservation,
the Hecke transform
has a unique coordinate vector $x'\in L^d$. Consequently
$T_{s,k}\Phi(x)=\Phi(x')$ for the precise basis map $\Phi$ of MR2--MR3.
This proves the rational clause without any assertion that its
matrix is integral at the residue characteristic of $s$.

\paragraph{MR7: prime-local Hecke integrality and one common denominator.}
Let $\mathfrak w$ be a prime of $\mathcal O_L$ whose residue
characteristic differs from that of $s$. Each representative in
(MR.10), together with its inverse, has denominators supported at
the rational prime below $s$. This follows directly for the
diagonal Hecke element, and multiplication by integral level
matrices preserves it. Thus the lattices attached to $g$ and
$ga_t$ are commensurable by an integer invertible at
$\mathfrak w$. They have equal localizations there.

The same is true of their actual $\nabla$ coefficient lattices.
To justify this last step, choose one basis of each localized
underlying lattice. The scalar-extended lattice, its dual, and all
polynomial coefficient modules defined from it use that same basis.
Their presentations therefore agree over $\mathcal O_{L,\mathfrak w}$.
This argument uses the actual dual-Weyl lattice $\nabla$; it makes
no assertion that integral $\Delta$ and $\nabla$ are equal.
Since $q_s$ is a unit at $\mathfrak w$, (MR.10) preserves the
localized coefficient lattice. If $x\in\mathcal N_{\psi}$ and $x'$
is its rational Hecke transform, each finite class $c$ admits
$d_c\in\mathcal O_L\setminus\mathfrak w$ with
$d_c(E_\psi x')_c\in\nabla(\Lambda_c)$. Put $d=\prod_c d_c$.
Then $d\notin\mathfrak w$ and every coordinate of $dx'$ is integral:
\[
                    dx'\in\mathcal N_\psi.          \tag{MR.11}
\]
This proves the exact denominator assertion of
\texttt{heckeOperator\_keeps\_rational\_integral\_levelForms}.

\paragraph{MR8: extending the fixed rational basis to the full MtReal map.}
Now use the actual tuple comparison
$e:L^d\simeq D.\mathrm{MtL}(\mu,j)$, and suppose, as in
\texttt{TupIn}, the vectors
$b_a=j_{\rm Mt}(e(e_a))$ are complex-independent and span the fixed
$D.\mathrm{Mt}(\mu,j)$. They are a basis. Define
\[
 \mathrm{real}\Bigl(\sum_a z_a b_a\Bigr)
             =\sum_a z_a\Phi(e_a).                  \tag{MR.12}
\]
MR2 proves injectivity and MR3 proves that the range is the entire
specified type space. Semilinearity gives
$\mathrm{real}(j_{\rm Mt}(e(x)))=\Phi(x)$ for every $x\in L^d$.
MR4 supplies the open level and MR5 the actual algebraic weight.
These are all seven non-instance fields of \texttt{MtReal}:
weight, level, character, openness, map, injectivity and full range.
This construction proves \texttt{side\_read} and \texttt{tup\_read}
under their literal geometric/basis hypotheses. Their exact maps
are specified by (MR.12); dimension counting is not used to invent
a map on an unrelated space.

\paragraph{MR9: all eight Good clauses and the exact MtInput lattice.}
The product, parahoric, unipotent and hyperspecial clauses are MR4;
the weight clause is MR5; rational Hecke stability is MR6.
For integrality take a generator $y$ of the fixed
$\mathrm{MtLat}(\mu,j)$. The given \texttt{OLatTie} writes
$y=j_{\rm Mt}(e(x))$ with $x\in\mathcal N_{\psi}$.
Use MR7, then write the integral vector $dx'$ in the given finite
generators as
\[
 j_{\rm Mt}(e(dx'))=\sum_g\epsilon(c_g)g,
       \qquad c_g\in\mathcal O_L.
\]
Because $d\notin\mathfrak w$, each $d^{-1}c_g$ has nonnegative
additive valuation at the specified place, and
\[
      T_{s,k}\mathrm{real}(y)
        =\mathrm{real}\Bigl(\sum_g\epsilon(d^{-1}c_g)g\Bigr).
                                                               \tag{MR.13}
\]
The generators already lie in the range of $j_{\rm Mt}$, and their
complex span is full by the same lattice tie. This is the eighth
\texttt{Good.full} clause. The exclusion of the residue prime in
MR7 is exactly the exclusion in \texttt{Good.int}.

Choose these realizations for every tuple and every $1\le j\le m$.
The result is \texttt{MtFam}. The source's generator-place condition
\texttt{GenAt}, together with the fixed conductor support and half
selection, places every Hecke generator outside both the level and
the relevant coefficient residue prime. Applying (MR.13) gives all
eight fields of \texttt{MtInput}, including the literal exponent
range $k\le r$ and the twisted-place qualification of its integral
and full-lattice clauses.

\paragraph{MR10: pointwise MtAd on used and unused tuples.}
Let $Y$ be the supplied \texttt{R3.Read} on the same fixed state.
For a used tuple its \texttt{eMt}, \texttt{formsL} and
\texttt{latL} are exactly the basis and lattice data in MR8--MR9.
The reading equality first holds on these basis vectors, by the
definition of \texttt{Y.mtRead} and the geometric coefficient map.
Evaluation at any fixed $g$ is complex-linear. Equality on the
basis therefore proves, for every $x$ and every $g$,
\[
 e_W\bigl(\mathrm{real}(x)(g)\bigr)
       =(Y.\mathrm{IL}.j_{\rm Mt}(\mu,h,i)(x))(P_i\psi(g)).
                                                               \tag{MR.14}
\]
Here $P_i\psi$ denotes the actual frame map $(PV_i).\psi$.
Its standardness, rational-point equivalence, coefficient-action
intertwining and entry-by-entry identity-frame equality follow from
the matrix construction in MR3 and MR5. These, (MR.14), and $e_W$
are precisely the seven \texttt{MtAd} fields.

Unused tuples cannot simply be omitted from \texttt{MtFam}.
The source passes to the all-tuples level, whose level is constructed
from the same half selection. The literal matrix equality
$m_L(\mu,j)=\mathrm{lMat}(\mu,j)$ transports the already supplied
\texttt{Rw.Tie1}: its complex basis, full span and integral-lattice
identity supply MR8 again. No reading equality is requested for an
unused tuple. Combining the used and unused constructions proves
\texttt{inMt\_state} and \texttt{sMt} on an actual supplied
\texttt{St}/\texttt{ColPin}. It does not construct that state or
the missing same-object ties required to obtain it.

\paragraph{MR11: what the rational input really supplies.}
For a rational form datum, a finite full lattice $M_{K_s}$ in
$W_{K_s}$ gives the literal \texttt{RatIn.fd} and
\texttt{RatIn.mk\_fg}. Choose a finite set of module generators
$t_1,\ldots,t_a$. Each lies in $M_{K_s}$, every lattice vector is an
$\mathcal O_{K_s}$-linear combination of them, and every rational
vector is a $K_s$-linear combination. A basis of $W_{K_s}$ which
remains a basis over $L$ gives \texttt{RatIn.ks}. The finite group
action and its character descend by their actual coefficient
identities; this supplies \texttt{rho} and \texttt{chi}.

The remaining assertions are substantial. With $H_{k,P}$ the
actual rational Hecke operators, $e$ the prescribed central
idempotent, and $U$ the prescribed finite unipotent group, put
\[
 W^{\rm good}=\ker(e-1)\cap
      \bigcap_{u\in U}\ker(a(u)-\psi_U(u)),           \tag{MR.15}
\]
where $a$ is the character-adjusted action. One must supply the
literal good-row decomposition
\[
 W^{\rm good}\subseteq
  \sum_{\substack{\sigma:\ \mathrm{isot}(j,\sigma)\ne0\\
                              \mathrm{goodRow}(\sigma)}}
    \left(W^{\rm good}\cap
        \bigcap_{\substack{1\le k\le r\\P\ {\rm split},\ P\notin S_b}}
            \ker(H_{k,P}-e_{j,\sigma,k,P})\right),   \tag{MR.16}
\]
with $e_{j,\sigma,k,P}\in K$, and the power-denominator bound
\[
    (\mathrm N P)^a H_{k,P}(M_L\cap W^{\rm good})
                       \subseteq M_L               \tag{MR.17}
\]
for some $a$, where $M_L=\mathcal O_L M_{K_s}$.
Together with the rational stability of $H$, $H_{0,P}=1$, the
specified nontrivial $\theta$ and the actual Hecke/eigenvalue
readings, these construct the fourteen fields of \texttt{RatIn}
and the five fields of \texttt{RatAd}. This is exactly the
\texttt{sRat\_of\_isotInternal} assembly. Its source derives
(MR.16)--(MR.17) from genuine internal row decomposition and
Hecke-transform/reading/integrality inputs. They are retained here;
finite class dimension alone does not prove either assertion.

\paragraph{MR12: an integral additive reading without enlarging the DVR.}
Let $O$ be a DVR, $H=\operatorname{Frac}(O)$, and
$\iota:K\hookrightarrow H$ the specified reading. There exists an
additive map $\lambda:L\to H$ such that
\[
 \lambda(1)=1,\qquad \lambda(ax)=\iota(a)\lambda(x)
 \ (a\in K),\qquad \lambda(\mathcal O_L)\subset O.    \tag{MR.18}
\]
Here is a direct construction. Extend $\iota$ to an embedding
$L\hookrightarrow\overline H$, and let $N$ be the $O$-span of the
images of an integral $\mathbb Z$-basis of $\mathcal O_L$. It is a
finite torsion-free $O$-module and contains $1$. Every element of $N$
is integral over $O$. Since $O$ is integrally closed,
\[
                      N\cap H\cdot1=O\cdot1.        \tag{MR.19}
\]
It follows that $N/O\cdot1$ is torsion-free: if a nonzero element
of $O$ annihilates a class, its representative lies in $H\cdot1$,
and (MR.19) applies. A finite torsion-free module over a DVR is free.
Lift a basis of this quotient, and define an $O$-linear functional
$N\to O$ to be $1$ on $1$ and $0$ on those lifts. Extend first
$H$-linearly to $HN$ and then to $\overline H$ by a vector-space
complement. Compose with $L\hookrightarrow\overline H$.
The resulting map has all three properties in (MR.18).

It is only additive and $K$-semilinear; multiplicativity is neither
claimed nor needed. This is why the argument works even when the
whole field $L$ does not embed into the already fixed $H$.
The source's \texttt{exists\_integral\_functional} also proves the
underlying primitive-vector splitting without assuming a full
lattice, by taking a basis of the image in $V/Hv$; the same proof
above applies to that image and extends the resulting functional.

\paragraph{MR13: finite integral equations descend to the original field.}
Choose a $K_s$-basis $b_1,\ldots,b_d$ of $W_{K_s}$ which is also
an $L$-basis of $W_L$, and define
\[
 \Phi\Bigl(\sum_a x_a b_a\Bigr)
           =\sum_a\lambda(x_a)(1\otimes b_a)
       \in H\otimes_{K_s}W_{K_s}.                  \tag{MR.20}
\]
Then $\Phi$ is $K$-semilinear, fixes rational vectors in the sense
$\Phi(w)=1\otimes w$, and commutes with every $L$-linear operator
whose restriction preserves $W_{K_s}$. This follows by writing its
matrix in the basis $b_a$: all its entries lie in $K_s\subset K$.
Also $\Phi(aw)=\lambda(a)(1\otimes w)$ for
$a\in L$, $w\in W_{K_s}$, so integral $a$ gives an $O$-multiple.

The further assertion needed for lattice stability is integral,
not merely vector-space descent. Suppose finitely many vectors
$v_{t,a}$ have all their $b$-coordinates in $K$, and
$o_t\in O$ satisfy
\[
                    \sum_t o_t\Phi(v_{t,a})=0
                                      \quad\text{for every }a.
                                                               \tag{MR.21}
\]
Clear one common denominator in the finitely many coordinates,
giving a matrix $A$ over $\mathcal O_K$. The embedding
$\mathcal O_K\to O$ is well-defined because images of global
integers are integral over the integrally closed ring $O$.
As an $\mathcal O_K$-module $O$ is torsion-free; since
$\mathcal O_K$ is Dedekind it is flat. For clarity, this flatness
can be checked locally: a localization of a Dedekind domain at a
nonzero prime is a DVR, and a torsion-free module over a DVR is a
filtered union of finite free modules, hence flat. Flatness is local.
Tensor the exact sequence
$0\to\ker A\to\mathcal O_K^t\to\mathcal O_K^a$ with $O$.
Thus every vector $(o_t)$ in (MR.21) is an $O$-linear combination
of images of solutions $(c_t)\in\mathcal O_K^t$ satisfying
$\sum_t c_tv_{t,a}=0$. Cancelling the nonzero common denominator
recovers the original equations. This proves the exact
finite-equation assertion of \texttt{exists\_integral\_descent}.

\paragraph{MR14: the full local lattice and its genuine good intersection.}
Fix a twist index $i$ from the already constructed \texttt{RingSpec}
and a side $b$. Put $O=O_i$, $H=\operatorname{Frac}(O_i)$, with
the side-dependent $K_s$-algebra structure obtained from
$\mathrm{read}_{i,b}$. For every actual row $1\le j\le m$, define
\[
 \mathcal F=H\otimes_{K_s}W_{K_s},\qquad
 B^{\rm sup}=O\langle1\otimes M_{K_s}\rangle.        \tag{MR.22}
\]
If $t_1,\ldots,t_a$ generate $M_{K_s}$, their images generate
$B^{\rm sup}$, since all images of $\mathcal O_{K_s}$ belong to
$O$. They span $\mathcal F$ over $H$. Consequently $B^{\rm sup}$
is finite, torsion-free and free over $O$.

Let $V_0\subset\mathcal F$ be the actual central-idempotent and
generic-character subspace, and put
\[
                       B=B^{\rm sup}\cap V_0.       \tag{MR.23}
\]
It is a submodule of a finite module over a noetherian ring, hence
finite, and is torsion-free, hence free. Moreover $HB=V_0$.
Indeed, express $v\in V_0$ in a finite generating set of
$B^{\rm sup}$ and clear the finitely many denominators to obtain
$0\ne d\in O$ with $dv\in B^{\rm sup}\cap V_0$.
The natural map $H\otimes_O B\to\mathcal F$ is injective:
a localized vector is zero only when a nonzero denominator kills
its numerator, impossible in a vector space. Its image is exactly
$V_0$. Thus (MR.23) is the source's \texttt{BI}, with its actual
rank, not a free module selected solely from that rank.

\paragraph{MR15: the actual idempotent, generic character and good space.}
The specified coefficients $e(g)\in K$ map under $\epsilon$ to
the central good-type projector $e_0$ in $\mathbb C[G_0]$.
Coefficientwise injectivity of $K[G_0]\to\mathbb C[G_0]$ reflects
$e_0^2=e_0$ and centrality. Thus $e^2=e$ and $e$ is central over
$K$, and the same polynomial identities hold after every actual
$\mathrm{read}_{i,b}$. The character-adjusted action is
\[
       a(g)=\chi(g)^{-1}\rho(g),\qquad
       a_{i,b}(g)=\mathrm{read}_{i,b}(\chi(g))^{-1}
                                      (1\otimes\rho(g)).\tag{MR.24}
\]
No ambiguity is introduced by changing the side. Each value of the
finite additive character $\theta_k$ has order dividing the prime
$g_k$, hence its complex conjugate is its inverse. The literal
conjugate-reading identity of \texttt{TwR} gives
$\mathrm{read}_{i,\mathrm{false}}(\zeta)
=\mathrm{read}_{i,\mathrm{true}}(\zeta)^{-1}$ for these values.
Injectivity preserves $\theta_k\ne1$. This proves the exact
\texttt{psiU\_side}, genericity and nonzero-coordinate conditions.

Let $p_U=|U|^{-1}\sum_u\psi_U(u)^{-1}a(u)$. The calculation of
MR2 makes it the projector onto the generic-character subspace;
centrality makes $e$ commute with it. Therefore
$V_0=e\operatorname{im}(p_U)$ is precisely the common kernel of
$e-1$ and $a(u)-\psi_U(u)$. These matrices are defined over $K$.
Finite-kernel descent implies that their common kernel over $H$
is spanned by the images under $\Phi$ of vectors in
$W^{\rm good}\subset W_L$. Apply (MR.16) to these vectors and
use the intertwining of MR13. This proves that $V_0$ is spanned by
the simultaneous eigenspaces with the actual row labels and
eigenvalues
\[
           c_{i,b,P}^{\,k}\,
                  \mathrm{read}_{i,b}(e_{j,\sigma,k,P}).\tag{MR.25}
\]
In passing between the $K$-kernel and the $L$-kernel, one may apply
the $K$-linear coefficient projection $\lambda$ to the $L$
coordinates; it is the identity on $K$. This proves the spanning
claim without extending $L\hookrightarrow H$ or changing $H$.

\paragraph{MR16: stability of the actual intersection lattice.}
Write $H'=c_{i,b,P}^{\,k}(1\otimes H_{k,P})$, where
$c_{i,b,P}\in O_i^\times$ is the prescribed delta correction.
Take $x\in B=B^{\rm sup}\cap V_0$ and write it as
$x=\sum_t o_t\Phi(w_t)$ with $w_t\in M_{K_s}$.
The equations defining $V_0$ say that $(o_t)$ satisfies a finite
system as in (MR.21): use the images of
$(e-1)w_t$ and $(a(u)-\psi_U(u))w_t$, whose coordinates lie in $K$.
MR13 expresses $x$ as an $O$-linear combination of $\Phi(y)$ with
\[
          y\in\mathcal O_L M_{K_s}\cap W^{\rm good}.\tag{MR.26}
\]
For such $y$, (MR.17) gives
$(\mathrm N P)^aH_{k,P}y\in\mathcal O_LM_{K_s}$.
MR12--MR13 send this lattice into $B^{\rm sup}$. Because
$P\notin S_0(i)$, it is not above the residue characteristic
$\ell_i$; thus $\mathrm N P$ is a unit of $O_i$. Cancelling its
power proves $\Phi(H_{k,P}y)\in B^{\rm sup}$.
The spectral expansion (MR.16) also proves
$H_{k,P}W^{\rm good}\subset W^{\rm good}$.
Hence $H'\Phi(y)\in V_0$, and multiplication by the unit
$c_{i,b,P}^{\,k}$ preserves both conditions. Extend over the
finite $O$-linear combination to obtain
\[
                              H'B\subseteq B.        \tag{MR.27}
\]
This is \texttt{HkStab} on the original coefficient ring. It is
proved for exactly $1\le j\le m$, $1\le k\le r$, split $P$
outside $S_0(i)$. The source defines the lattice operator to be
the identity outside this range, in particular at $k=0$; it does
not claim (MR.27) for a ramified or excluded Hecke operator.

\paragraph{MR17: integral eigenvalues and finite, separated signatures.}
An endomorphism of a finite free $O$-module satisfies its monic
characteristic polynomial. If its scalar extension has a nonzero
eigenvector of eigenvalue $a\in H$, applying that polynomial to
the eigenvector gives $p(a)=0$. Since $O$ is integrally closed,
$a\in O$. The same conclusion for a finite module follows from
the determinant trick; freeness from MR14 suffices here.

For fixed $i$, take triples $(j,b,\sigma)$ with
$1\le j\le m$, $\mathrm{isot}(j,\sigma)\ne0$ and
$\mathrm{goodRow}(\sigma)$ for which every admissible eigenvalue
in (MR.25) lies in $O_i$. Store that integral eigenvalue at the
admissible pairs $(k,P)$ and store $1$ elsewhere. Define
$\mathrm{Sig}(i)$ as the image of these triples under
\[
       (j,b,\sigma)\longmapsto
                (j,b,(k,P\mapsto h_{j,b,\sigma}(k,P))).\tag{MR.28}
\]
This is a finite set even though the eigenvalue function has
infinitely many arguments, because its domain of triples is finite.
Select an original row label for each element of this image.
Two signatures with the same row number, side and all admissible
eigenvalues are equal: outside that range both functions are $1$.
This proves the literal injectivity clause and $h(0,P)=1$.

MR15 spans $V_0$ by simultaneous eigenvectors. Every nonzero
simultaneous eigenvector used in that span has integral eigenvalues
by MR16 and the preceding
monic-polynomial argument. Its triple is therefore included in
(MR.28); zero vectors contribute nothing. Via MR14, these
signature eigenspaces span $H\otimes_O B$. The selected original
row is good and has nonzero prescribed isotypic space, and
(MR.25) is unchanged. These are all five conclusions of
\texttt{exists\_sigData\_hecke\_eigenvalue\_formula}.

\paragraph{MR18: characteristic polynomials and the positive TwF fields.}
For an admissible individual Hecke operator, MR17 supplies a basis
of each of its eigenspaces and their direct sum spans the fraction
space. Distinct eigenvalues have independent eigenspaces: applying
$\prod_{a\ne a_0}(H'-a)$ isolates each summand. The matrix in the
collected basis is diagonal, so its characteristic polynomial is
the product of the corresponding linear factors. The matrix of
the integral restriction in an $O$-basis maps to the matrix of
its scalar extension. Determinants commute with that map; hence
the literal integral characteristic polynomial splits over $H$.
For a non-admissible operator the defined matrix is the identity,
and the polynomial is $(X-1)^{\operatorname{rank}_O B}$.

This constructs the positive form part of \texttt{TwF}: the exact
finite signature set, its row number and side, the module $B$ of
(MR.23), its finite freeness, the restricted Hecke operators, and
the integral eigenvalues. The signature Galois representation is
the actual previously constructed $\rho_{i,\sigma,b}$ of NR24
at the selected original row and side. Its continuity and its
universal reading clause follow from NR24, under its explicit
\texttt{GalEx}, \texttt{RowCusp}, \texttt{ChebSplit} and
\texttt{SideA} premises. Signature selection does not weaken
those premises or change that representation.

\paragraph{MR19: explicit weights over the same fraction fields.}
The row parameter is genuinely realizable by a dominant integral
weight. Let $n=R.n_V$ and let $S=\operatorname{paramV}(R,h,j,s)$.
Its entries are the multiset
\[
 \{v_a\}_{a<k-1}\ \sqcup\ \{-b_a\}_{a<m}\ \sqcup\
 \{-e_j+(t-1-2c)/2\}_{c<t},
          \quad\text{then shifted by }-s.           \tag{MR.29}
\]
The source's adjacent gap conditions make each block strictly
decreasing and separate the three blocks: the third lies between
the appropriate consecutive $-b$ entries and below the $v$ entries.
Thus $S$ has $n$ distinct entries. Integrality of $v,b,e$, oddness
of $t$, and evenness of $n$ show that, when $s\in\mathbb Z+1/2$,
each entry is $z+(n-1)/2$ with $z\in\mathbb Z$.
List the negatives of these integers increasingly as
$a_0<\cdots<a_{n-1}$ and set
\[
                         \lambda_i=-a_i+i.          \tag{MR.30}
\]
Since $a_{i+1}\ge a_i+1$, $\lambda$ is weakly decreasing.
Its infinitesimal-character entries are
$\lambda_i+(n-1)/2-i=-a_i+(n-1)/2$, exactly $S$.
The actual character frame gives $s=1/2$ for the one-dimensional
$\eta_T$ parameter; more generally the proof only needs the stated
half-integrality. This proves the literal combinatorial weight
supplier rather than leaving it as a record field.

For every $i,j,b$, use the side-dependent embeddings
$E\to\operatorname{Frac}(O_i)$ prescribed by \texttt{RingSpec}
and take the dual of the tensor polynomial module $VV(\lambda)$.
Its finite basis and matrices are over that very field. The
highest-weight input of MR5 proves irreducibility and the exact
dual highest weights; extension of scalars preserves this tensor
polynomial realization and its dual, hence preserves
irreducibility. The comparison is explicit on the polynomial
basis, including the inverse determinant powers. Therefore the
weight bound $c_2$ in \texttt{wt\_weights} can be the bottom field:
no additional local field is needed for this weight construction.
Only the source's harmless out-of-row padding uses a default
weight; every row $1\le j\le m$ uses (MR.29)--(MR.30).

\paragraph{MR20: actual local levels and the whole coefficient-field type space.}
Fix an admissible character $\psi$ and conductor set $Q$ at
$(i,k)$. Select one prime in each conjugate pair using the prescribed
$P_q$; write $Q_P=\operatorname{qpOf}(P_q,Q)$ and
$D_P=\operatorname{qpOf}(P_q,p_D(i))$. The split hypotheses and
conjugation stability show that each pair is met once. The three
families of real places, those from $Q_P$, $D_P$ and the type
primes, are disjoint. At these places define $J$ inside the
integral maximal compact by, respectively, the $(r,n-r)$
parabolic, the same parabolic, and the type-prime unipotent level.

These are subgroups by reduction modulo the local maximal ideal:
block upper-triangular matrices form a subgroup, and requiring the
upper-left block to be upper unitriangular is also a subgroup.
The determinant of the complementary block is a unit and
multiplicative in the residue field. The superdiagonal sum of an
upper unitriangular block is additive under multiplication.
Consequently the following product is a character on $J$:
the actual generic additive/type determinant factor at every type
prime, the tame character of $\psi$ on each complementary
determinant in $Q_P$, and the tame character of
$\delta_i^{\varepsilon_b}$ on the corresponding determinants in
$D_P$, where $\varepsilon_{\rm true}=1$ and
$\varepsilon_{\rm false}=-1$. Tame residue values are well-defined
and multiplicative by GA17 and AR3; wild inertia is killed since
the character orders and these residue characteristics are
coprime. At the type primes the literal \texttt{tauT} is the
residue character of the already normalized \texttt{chi1}, not
an arbitrary finite-order replacement.

The product of the norms of these finitely many real places gives
a positive integer $N$ for which the congruence subgroup is in
$\ker\chi$. Outside those places $J$ is maximal compact with
trivial character. These prove all local, open and hyperspecial
clauses of \texttt{LevLawsT}, with the literal sets $Q_P,D_P$.

Finally, the whole space
$\operatorname{TypeSpF}(\rho_{i,j,b,k},J,\chi)$ is finite-dimensional.
Choose finitely many classes modulo the maximal compact, then
finitely many coset representatives of its level-$N$ congruence
subgroup. Evaluation at their products is injective: every group
element is rational times one such product times an element of
the congruence kernel, and equivariance recovers its value.
The target is a finite product of the finite-dimensional coefficient
space. The literal \texttt{MtSubT} chooses this actual type space
when these level laws exist, and its comparison map is the identity
under that equality. Take the $O_{i,k}$-span of any fraction-field
basis as \texttt{MtLatT}; it is finite free and spans the space.
Its comparison to the original global \texttt{MtLat} is a separate
arithmetic compatibility problem, not a consequence of this choice
of basis.

\paragraph{MR21: the initial NExpTied fields and exact normalized comparisons.}
MR14--MR20 construct the coefficient-field form modules, the
generic action, central idempotent, weight modules, and full
admissible type spaces used in the source's \texttt{mkEx}.
The central coefficients remain the original $e(g)\in K$.
The universal anticyclotomic characters and their support,
continuity and conjugation laws are NR8--NR11. Their inertia
surjectivity is only asserted under literal \texttt{AdmQ}:
that condition supplies, at each conductor pair, a character
surjective on the chosen inertia. Its conjugate is surjective on
the other prime by inversion. The universal character's coordinate
is the selected two-prime character, so it is surjective there and
zero on inertia at a different, nonconjugate conductor prime.
Thus no stronger inertia-surjectivity assumption has been inferred
from support alone. The source signs $\mathrm{sgnT}$ and
$\mathrm{sgnG}$ are both $1$; MR15 and MR20 check their actual
character equalities with this convention.

For admissible rows the two tensor descriptions
\texttt{FmT} and \texttt{Fm S} are identical after replacing
$\mathrm{WKr}(j)$ by $W_{K_s}$ and identifying their identical
$K_s$-algebra structures. Send $a\otimes w$ to $a\otimes w$.
This gives the literal $e_F$, carries $B^{\rm sup}$ to the
specified lattice, and intertwines the character-adjusted action.
The Hecke comparison is
\[
 e_F(H'_{k,P}x)=c_{i,b,P}^{\,k}
                         (1\otimes H_{k,P})e_F(x).  \tag{MR.31}
\]
The unit $c_{i,b,P}$ is the value of the already fixed delta
representation at the selected Frobenius, inverted on the false
side. Any two allowed Frobenius representatives differ by inertia,
on which delta is trivial outside $S_0(i)$, so the value is
independent of the selected representative. This is the literal
\texttt{cHk\_del}, with the existing \texttt{IsFrob} convention.
The lattice injection is the subtype map $B\hookrightarrow
\mathcal F$. Its Hecke equality is (MR.27) for $k>0$ and the
identity equality for $k=0$. All these comparisons use the original
\texttt{RingSpec} and the same two side embeddings.

\paragraph{MR22: exact SideB conclusion and its initial-stage boundary.}
Retain the full input list of \texttt{second\_step}: the fixed
$D$, character frame, $j_E$, \texttt{SideA}, actual rational datum
and \texttt{RatIn}, a finite low-prime set, $P_q$ lying above every
real place, an actual \texttt{MtFam}/\texttt{MtInput}, global Artin
reciprocity, \texttt{GalEx}, \texttt{RowCusp}, \texttt{ChebSplit},
conjugation stability of $K$, integrality of $e\mathrm{Exp}$, and
splitting at both auxiliary primes. MR8--MR10 can supply the
Mt-fields when their compatible geometric state is given; MR11
specifies exactly the remaining rational spectral inputs.

NR24 chooses a finite field bound $c_1(\ell,q)$ for all genuine
row representations before the later parameters. MR19 chooses
$c_2=\bot$. Put $c_0=c_1\vee c_2$. Thus the quantifier order is
\[
 \exists c_0\ \forall b_0,b_1,S_x,T,\mathrm{Sp},\qquad
 c_0\le b_0,\quad S_b\subseteq S_x,\quad
 \bar S_x=S_x,\quad
 \{P:P\mid p,\ p\in\mathrm{low}\}\subseteq S_x.    \tag{MR.32}
\]
For every such \texttt{RingSpec}, MR14--MR21 give the actual
form/signature fields and the initial \texttt{NExpTied}. The
remaining fields of \texttt{SideB} have explicit sources:
the coefficient frame, normalized characters and all embedding
squares are the same \texttt{Sp}; cover at every prescribed
coefficient place is NR17; rational action is the same
\texttt{RatIn}; $e_F$, the lattice and Hecke equalities are MR21;
signature span, separation and eigenvalues are MR17--MR18; and
\texttt{rhoS} is NR24 at the selected good row. This proves every
\texttt{SideB} field under the source's exact qualification below.

That qualification is material. Literal \texttt{mkF/mkEx} initialize
\texttt{Theta} and \texttt{ThetaPsi} as zero submodules,
\texttt{Dat=PEmpty}, \texttt{thP=0}, \texttt{twDel=0}, and
$e_H=1$; the toric family is \texttt{TwL.empty} and
\texttt{SideA.bare} says that the old local index set is empty.
These fields are initialization choices, not the genuine theta
image, toric realization or nonempty local family. The source
conclusion is exactly
\[
 \exists\mathrm{TwF},\mathrm{ex},h_0,\quad
  [\text{the full-lattice properties of }B^{\rm sup}]
  \ \land\ \forall\mathrm{th}:NTh(\mathrm{new\ data}),\quad
        \operatorname{Nonempty}(\mathrm{SideB}(D_{\rm new})).
                                                               \tag{MR.33}
\]
It does not assert that such a \texttt{th} has been constructed.
$D_{\rm new}$ is the specified \texttt{setTw} datum: it preserves
the original host, row spaces, $L$, $K$, normalized frame and
rational operators, and replaces its twist/form expansion fields
by exactly the objects above. Its twist index set can be the
genuinely nonempty all-prime family from NR14--NR18. Nonvanishing
of a good space is faithfully preserved by (MR.22)--(MR.23) and
MR15, but is not asserted for every row if the original space is
zero. No new nonempty local index, genuine theta operator or
\texttt{IsTripleFor} is supplied by (MR.33).

\paragraph{MR23: one finite local field splits the entire commuting Hecke family.}
Fix the original row $j$, coefficient place $w$, embedding
$\iota_w:L\hookrightarrow\overline{\mathbb Q}_\ell$, and actual
finite-dimensional space
$V_L=D.\mathrm{what}_2.\mathrm{MtL}(\mu_w,j)$.
Let $H_\alpha\in\operatorname{End}_L(V_L)$ be the source's
\texttt{hLC} family, indexed by all pairs of finite places and
Hecke indices, with its supplied pairwise commutation law.
The $L$-linear span of this possibly infinite family is finite
dimensional. Choose finitely many members $H_1,\ldots,H_s$
spanning it, and put
\[
             p(X)=\prod_{a=1}^{s}\det(X-H_a)\in L[X].
                                                               \tag{MR.34}
\]
This is a nonzero monic polynomial and annihilates every selected
$H_a$ by Cayley--Hamilton. The empty product is $1$ when the
operator span is zero. Adjoin to $\mathbb Q_\ell$ all roots of
$\iota_w(p)$ in its algebraic closure. There are finitely many
roots, each algebraic over $\mathbb Q_\ell$, so the resulting
$F_0/\mathbb Q_\ell$ is finite.

Now let $H$ be any later coefficient field whose embedding in
$\overline{\mathbb Q}_\ell$ contains $F_0$ and extends the same
$\iota_w$. Every root of the polynomial obtained from $p$ over
$H$ therefore lies in $H$, so each selected operator has split
minimal polynomial on $H\otimes_LV_L$. Commuting split operators
are simultaneously triangularizable over $H$: a nonzero
eigenspace of the first operator is invariant under the others;
successively taking nonzero eigenspaces gives a common
eigenvector, and induction on the dimension applied to the
quotient gives an invariant full flag. Restrictions and quotients
remain split because their minimal polynomials divide the
original ones. Every $H$-linear combination of the selected
operators preserves this flag. Hence every original $H_\alpha$
has split characteristic and minimal polynomials over $H$.

For a split minimal polynomial
$\prod_c(X-c)^{a_c}$, the factors for distinct $c$ are coprime.
B\'ezout polynomials give the direct-sum decomposition into
$\ker(H_\alpha-c)^{a_c}$. Consequently
\[
 \sum_{c\in H}\operatorname{GenEig}(1\otimes H_\alpha,c)
          =H\otimes_LV_L.                              \tag{MR.35}
\]
This is the literal \texttt{exists\_finExt\_split} conclusion
for all admissible $\alpha$, with the stronger elementary
argument valid for every member of the family. No semisimplicity
has been assumed or concluded. The quantifier is one $F_0$ for
the fixed $(j,w)$, before every later compatible \texttt{CoefIn};
it is not one finite global field for all coefficient places.

\paragraph{MR24: a nonempty residual Hecke block has a genuinely local algebra.}
Let $O$ be a local ring with maximal ideal $\mathfrak m$, let
$X$ be a finite nonempty set, and let
$T\subseteq O^X$ be an $O$-subalgebra whose coordinate values
are congruent modulo $\mathfrak m$:
\[
             f(x)-f(y)\in\mathfrak m
       \quad(f\in T,\ x,y\in X).                       \tag{MR.36}
\]
Fix $x_0\in X$. If $f(x_0)$ is a unit, (MR.36) makes every
$f(x)$ a unit. The polynomial
$P_f(Z)=\prod_{x\in X}(Z-f(x))$ annihilates $f$ in $O^X$;
its coefficients belong to $O$, hence are constants in $T$,
and its constant coefficient is a unit. Writing
$P_f(Z)=a_0+ZQ(Z)$ gives
$f[-a_0^{-1}Q(f)]=1$ in $T$. Thus $f$ is a unit of $T$
if and only if $f(x_0)$ is a unit of $O$. The inverse implication
also follows directly by applying coordinate evaluation.
The nonunits are therefore the kernel of the surjective map
\[
               T\longrightarrow O/\mathfrak m,
                   \qquad f\longmapsto f(x_0)\bmod\mathfrak m.
                                                               \tag{MR.37}
\]
Surjectivity follows from the constant functions. This kernel
is the unique maximal ideal, proving that $T$ is local and that
its residue field is the original $O/\mathfrak m$.

For the source's residual block, points are $O$-algebra
characters $\psi$ of the same Hecke algebra with a common
kernel after residue reduction. If $\psi,\phi$ are two such
characters, apply equality of their residue kernels to
$a-\phi(a)\cdot1$. Its $\phi$-value is zero, so
$\psi(a)-\phi(a)\in\mathfrak m$. This proves (MR.36) on the
Hecke generators; sums, products and constants preserve it,
so it holds on their entire generated algebra. The actual
\texttt{HeckeStd0}, finite-dimensionality and existing
\texttt{LocT} hypotheses give the finite nonempty point set
\texttt{XOf}. Applying the preceding argument to its actual
\texttt{TOf} proves the literal local-ring assertion. This is
locality of an existing nonzero-support block, not an existence
claim for a block at an arbitrary place.

\paragraph{MR25: every existing supported Hecke point has a nonzero eigenspace.}
Let $V$ be finite dimensional over a field $H$, let $S$ be a
commuting family of endomorphisms, and let $0\ne W\subseteq V$
be invariant under $S$. Suppose for every $a\in S$ that
$W$ lies in the generalized $\chi(a)$-eigenspace of $a$.
Choose a nonzero invariant subspace $U\subseteq W$ of smallest
positive dimension. For fixed $a\in S$ and $0\ne u\in U$,
choose the least $r>0$ with $(a-\chi(a))^ru=0$.
Then $(a-\chi(a))^{r-1}u$ is a nonzero element of
$U\cap\ker(a-\chi(a))$. This intersection is invariant under
all of $S$, because the operators commute. Its dimension cannot
be smaller than that of $U$, so it equals $U$. The argument is
valid for every $a$, without assuming that $S$ is finite.
Every nonzero $u\in U$ is therefore a simultaneous eigenvector.

In the actual localization construction, a point
$x\in\texttt{XOf}(i)$ includes the nonzero generalized
character space $W_x$. Apply the preceding argument to
$W_x$ and the admissible normalized Hecke operators. It gives
$u\ne0$ with their exact point values. Retain the actual
\texttt{MtLoc}, \texttt{ActEx}, and \texttt{ActHK} comparison:
its injective map $\Psi_i$ has range the prescribed residual
block, and it intertwines the specified Hecke action on the
localized relation lattice with those operators. Thus
$u=\Psi_i(z)$ for some $z\ne0$, and injectivity gives
\[
       a\,z=\operatorname{ev}_x(a)z
              \quad\text{for each Hecke generator }a. \tag{MR.38}
\]
The equality propagates to constants, sums and products, hence
to the entire generated algebra $T_i$. Therefore its actual
character eigenspace contains $z\ne0$; this is
\texttt{heckeEig\_ne\_bot}. The comparison $\Psi_i$ and
the point's nonzero support are essential retained premises.
Neither is inferred merely from MR23's generalized splitting
or from MR24's local-ring structure.

\paragraph{MR26: the original relation lattice stays free over the chosen DVR.}
Let $A$ be the original integral coefficient ring. The literal
\texttt{OsConcrete} condition says its given injection into
$L$ has image
$\{x\in F:x\text{ is integral over }\mathbb Z\}$ for a
number subfield $F\subseteq L$. Thus it identifies $A$ with
$\mathcal O_F$, so $A$ is Dedekind. It does not permit
replacement by an unrelated coefficient ring.

Let $v_1,\ldots,v_d$ be the specified complex generators of
the original lattice and let
\[
 N=\{(a_1,\ldots,a_d)\in A^d:
                 \sum_t\epsilon(a_t)v_t=0\}.          \tag{MR.39}
\]
If $0\ne c\in A$ and $ca\in N$, then
$\epsilon(c)\sum_t\epsilon(a_t)v_t=0$; the fixed coefficient
embedding is injective, so $a\in N$. Consequently $A^d/N$
is torsion free. It is finite, hence projective over the Dedekind
ring $A$: its localization at each nonzero prime is a finite
torsion-free module over a DVR and is free; finite presentation
and these local bases give projectivity. One may see the last
step by choosing finitely many denominator neighborhoods on
which bases exist and applying the finite locally free
criterion. The quotient map therefore has a section, giving
an $A$-linear projection $e:A^d\to N$ with $e|_N=1$.

For the prescribed ring homomorphism $s:A\to O$ to the chosen
DVR, apply $s$ to the matrix of $e$, obtaining $e_O:O^d\to O^d$.
Its image is exactly
$N_O=O\langle s(N)\rangle$: it maps into this span because
its columns are images of members of $N$, and it fixes the
span because it fixes every $s(n)$. Thus $e_O$ is an
idempotent and
\[
                 O^d/N_O\simeq\ker e_O.              \tag{MR.40}
\]
The right side is a finite torsion-free $O$-module, so it is
free. This proves \texttt{bjC\_free} for the literal quotient
of the original generator relations, for every existing
localization index and its actual coefficient map. It does not
identify that quotient with a different lattice by a dimension
comparison.

The remaining nonempty-realization problem now has a sharper
boundary. MR23--MR26 supply the finite splitting-field, local
block, nonzero-eigenvector and relation-freeness primitives.
The full \texttt{NdSpec}/\texttt{LocAll} constructor still
requires the compatible original \texttt{NChoice},
\texttt{ChoiceSide}, all coefficient/twist embedding squares,
the actual $\Psi$ comparison and local Galois/automorphic
specification. The initial object in MR22 does not supply these
requirements, nor a compatible nonplaceholder theta witness.


